%%%%%%%%%%%%%%%%%%%%%%%%%%%%%%%%%%%%%%%%%%%%%%%%%%%%%%%%%%%%%%%%%%%%%%%%%%%%%%%%
% The Legrand Orange Book
% LaTeX Template
% Version 2.2 (30/3/17)
%
% This template has been downloaded from:
% http://www.LaTeXTemplates.com
%
% Original author:
% Mathias Legrand (legrand.mathias@gmail.com) with modifications by:
% Vel (vel@latextemplates.com)
%
% License:
% CC BY-NC-SA 3.0 (http://creativecommons.org/licenses/by-nc-sa/3.0/)
%
% Compiling this template:
% This template uses biber for its bibliography and makeindex for its index.
% When you first open the template, compile it from the command line with the
% commands below to make sure your LaTeX distribution is configured correctly:
%
% 1) pdflatex main
% 2) makeindex main.idx -s StyleInd.ist
% 3) biber main
% 4) pdflatex main
%
% After this, when you wish to update the bibliography/index use the appropriate
% command above and make sure to compile with pdflatex several times
% afterwards to propagate your changes to the document.
%
% This template also uses a number of packages which may need to be
% updated to the newest versions for the template to compile. It is strongly
% recommended you update your LaTeX distribution if you have any
% compilation errors.
%
% Important note:
% Chapter heading images should have a 2:1 width:height ratio,
% e.g. 920px width and 460px height.
%
%%%%%%%%%%%%%%%%%%%%%%%%%%%%%%%%%%%%%%%%%%%%%%%%%%%%%%%%%%%%%%%%%%%%%%%%%%%%%%%%

%----------------------------------------------------------------------------------------
%	DOCUMENT CONFIGURATIONS
%----------------------------------------------------------------------------------------

\documentclass[12pt,a4paper]{book} % Default font size and left-justified equations

%%%%%%%%%%%%%%%%%%%%%%%%%%%%%%%%%%%%%%%%%
% The Legrand Orange Book
% Structural Definitions File
% Version 2.0 (9/2/15)
%
% Original author:
% Mathias Legrand (legrand.mathias@gmail.com) with modifications by:
% Vel (vel@latextemplates.com)
%
% This file has been downloaded from:
% http://www.LaTeXTemplates.com
%
% License:
% CC BY-NC-SA 3.0 (http://creativecommons.org/licenses/by-nc-sa/3.0/)
%
%%%%%%%%%%%%%%%%%%%%%%%%%%%%%%%%%%%%%%%%%

%----------------------------------------------------------------------------------------
%	VARIOUS REQUIRED PACKAGES AND CONFIGURATIONS
%----------------------------------------------------------------------------------------

\usepackage[top=3cm,bottom=3cm,left=3cm,right=3cm,headsep=10pt,a4paper]{geometry} % Page margins

\usepackage{graphicx} % Required for including pictures
%\graphicspath{{Pictures/}} % Specifies the directory where pictures are stored % uncomment to set pictures directory

\usepackage{float} % Required to fix floating objects

\usepackage{lipsum} % Inserts dummy text

\usepackage{tikz} % Required for drawing custom shapes

\usepackage[french]{babel} % French language/hyphenation

\usepackage{enumitem} % Customize lists
\setlist{nolistsep} % Reduce spacing between bullet points and numbered lists

\usepackage{booktabs} % Required for nicer horizontal rules in tables

\usepackage{xcolor} % Required for specifying colors by name
\definecolor{ocre}{RGB}{243,102,25} % Define the orange color used for highlighting throughout the book

\usepackage[outputdir=build]{minted}

%----------------------------------------------------------------------------------------
%	FONTS
%----------------------------------------------------------------------------------------

\usepackage{avant} % Use the Avantgarde font for headings
%\usepackage{times} % Use the Times font for headings
\usepackage{mathptmx} % Use the Adobe Times Roman as the default text font together with math symbols from the Sym­bol, Chancery and Com­puter Modern fonts

\usepackage{microtype} % Slightly tweak font spacing for aesthetics
\usepackage[utf8]{inputenc} % Required for including letters with accents
\usepackage[T1]{fontenc} % Use 8-bit encoding that has 256 glyphs

%----------------------------------------------------------------------------------------
%	BIBLIOGRAPHY AND INDEX
%----------------------------------------------------------------------------------------

\usepackage[style=alphabetic,citestyle=numeric,sorting=nyt,sortcites=true,autopunct=true,babel=hyphen,hyperref=true,abbreviate=false,backref=true,backend=biber]{biblatex}
\addbibresource{bibliography.bib} % BibTeX bibliography file
\defbibheading{bibempty}{}

\usepackage{calc} % For simpler calculation - used for spacing the index letter headings correctly
\usepackage{makeidx} % Required to make an index
\makeindex % Tells LaTeX to create the files required for indexing

%----------------------------------------------------------------------------------------
%	MAIN TABLE OF CONTENTS
%----------------------------------------------------------------------------------------

\usepackage{titletoc} % Required for manipulating the table of contents

\contentsmargin{0cm} % Removes the default margin

% Part text styling
\titlecontents{part}[0cm]
{\addvspace{20pt}\centering\large\bfseries}
{}
{}
{}

% Chapter text styling
\titlecontents{chapter}[1.25cm] % Indentation
{\addvspace{12pt}\large\sffamily\bfseries} % Spacing and font options for chapters
{\color{ocre!60}\contentslabel[\Large\thecontentslabel]{1.25cm}\color{ocre}} % Chapter number
{\color{ocre}}
{\color{ocre!60}\normalsize\;\titlerule*[.5pc]{.}\;\thecontentspage} % Page number

% Section text styling
\titlecontents{section}[1.25cm] % Indentation
{\addvspace{3pt}\sffamily\bfseries} % Spacing and font options for sections
{\contentslabel[\thecontentslabel]{1.25cm}} % Section number
{}
{\hfill\color{black}\thecontentspage} % Page number
[]

% Subsection text styling
\titlecontents{subsection}[1.25cm] % Indentation
{\addvspace{1pt}\sffamily\small} % Spacing and font options for subsections
{\contentslabel[\thecontentslabel]{1.25cm}} % Subsection number
{}
{\ \titlerule*[.5pc]{.}\;\thecontentspage} % Page number
[]

% List of figures
\titlecontents{figure}[0em]
{\addvspace{-5pt}\sffamily}
{\thecontentslabel\hspace*{1em}}
{}
{\ \titlerule*[.5pc]{.}\;\thecontentspage}
[]

% List of tables
\titlecontents{table}[0em]
{\addvspace{-5pt}\sffamily}
{\thecontentslabel\hspace*{1em}}
{}
{\ \titlerule*[.5pc]{.}\;\thecontentspage}
[]

%----------------------------------------------------------------------------------------
%	MINI TABLE OF CONTENTS IN PART HEADS
%----------------------------------------------------------------------------------------

% Chapter text styling
\titlecontents{lchapter}[0em] % Indenting
{\addvspace{15pt}\large\sffamily\bfseries} % Spacing and font options for chapters
{\color{ocre}\contentslabel[\Large\thecontentslabel]{1.25cm}\color{ocre}} % Chapter number
{}
{\color{ocre}\normalsize\sffamily\bfseries\;\titlerule*[.5pc]{.}\;\thecontentspage} % Page number

% Section text styling
\titlecontents{lsection}[0em] % Indenting
{\sffamily\small} % Spacing and font options for sections
{\contentslabel[\thecontentslabel]{1.25cm}} % Section number
{}
{}

% Subsection text styling
\titlecontents{lsubsection}[.5em] % Indentation
{\normalfont\footnotesize\sffamily} % Font settings
{}
{}
{}

%----------------------------------------------------------------------------------------
%	PAGE HEADERS
%----------------------------------------------------------------------------------------

\usepackage{fancyhdr} % Required for header and footer configuration

\pagestyle{fancy}
\renewcommand{\chaptermark}[1]{\markboth{\sffamily\normalsize\bfseries\chaptername\ \thechapter.\ #1}{}} % Chapter text font settings
\renewcommand{\sectionmark}[1]{\markright{\sffamily\normalsize\thesection\hspace{5pt}#1}{}} % Section text font settings
\fancyhf{} \fancyhead[LE,RO]{\sffamily\normalsize\thepage} % Font setting for the page number in the header
\fancyhead[LO]{\rightmark} % Print the nearest section name on the left side of odd pages
\fancyhead[RE]{\leftmark} % Print the current chapter name on the right side of even pages
\renewcommand{\headrulewidth}{0.5pt} % Width of the rule under the header
\addtolength{\headheight}{2.5pt} % Increase the spacing around the header slightly
\renewcommand{\footrulewidth}{0pt} % Removes the rule in the footer
\fancypagestyle{plain}{\fancyhead{}\renewcommand{\headrulewidth}{0pt}} % Style for when a plain pagestyle is specified

% Removes the header from odd empty pages at the end of chapters
\makeatletter
\renewcommand{\cleardoublepage}{
\clearpage\ifodd\c@page\else
\hbox{}
\vspace*{\fill}
\thispagestyle{empty}
\newpage
\fi}

%----------------------------------------------------------------------------------------
%	THEOREM STYLES
%----------------------------------------------------------------------------------------

\usepackage{amsmath,amsfonts,amssymb,amsthm} % For math equations, theorems, symbols, etc

\newcommand{\intoo}[2]{\mathopen{]}#1\,;#2\mathclose{[}}
\newcommand{\ud}{\mathop{\mathrm{{}d}}\mathopen{}}
\newcommand{\intff}[2]{\mathopen{[}#1\,;#2\mathclose{]}}
\newtheorem{notation}{Notation}[chapter]

% Boxed/framed environments
\newtheoremstyle{ocrenumbox}% % Theorem style name
{0pt}% Space above
{0pt}% Space below
{\normalfont}% % Body font
{}% Indent amount
{\small\bf\sffamily\color{ocre}}% % Theorem head font
{\;}% Punctuation after theorem head
{0.25em}% Space after theorem head
{\small\sffamily\color{ocre}\thmname{#1}\nobreakspace\thmnumber{\@ifnotempty{#1}{}\@upn{#2}}% Theorem text (e.g. Theorem 2.1)
\thmnote{\nobreakspace\the\thm@notefont\sffamily\bfseries\color{black}---\nobreakspace#3.}} % Optional theorem note
\renewcommand{\qedsymbol}{$\blacksquare$}% Optional qed square

\newtheoremstyle{blacknumex}% Theorem style name
{5pt}% Space above
{5pt}% Space below
{\normalfont}% Body font
{} % Indent amount
{\small\bf\sffamily}% Theorem head font
{\;}% Punctuation after theorem head
{0.25em}% Space after theorem head
{\small\sffamily{\tiny\ensuremath{\blacksquare}}\nobreakspace\thmname{#1}\nobreakspace\thmnumber{\@ifnotempty{#1}{}\@upn{#2}}% Theorem text (e.g. Theorem 2.1)
\thmnote{\nobreakspace\the\thm@notefont\sffamily\bfseries---\nobreakspace#3.}}% Optional theorem note

\newtheoremstyle{blacknumbox} % Theorem style name
{0pt}% Space above
{0pt}% Space below
{\normalfont}% Body font
{}% Indent amount
{\small\bf\sffamily}% Theorem head font
{\;}% Punctuation after theorem head
{0.25em}% Space after theorem head
{\small\sffamily\thmname{#1}\nobreakspace\thmnumber{\@ifnotempty{#1}{}\@upn{#2}}% Theorem text (e.g. Theorem 2.1)
\thmnote{\nobreakspace\the\thm@notefont\sffamily\bfseries---\nobreakspace#3.}}% Optional theorem note

% Non-boxed/non-framed environments
\newtheoremstyle{ocrenum}% % Theorem style name
{5pt}% Space above
{5pt}% Space below
{\normalfont}% % Body font
{}% Indent amount
{\small\bf\sffamily\color{ocre}}% % Theorem head font
{\;}% Punctuation after theorem head
{0.25em}% Space after theorem head
{\small\sffamily\color{ocre}\thmname{#1}\nobreakspace\thmnumber{\@ifnotempty{#1}{}\@upn{#2}}% Theorem text (e.g. Theorem 2.1)
\thmnote{\nobreakspace\the\thm@notefont\sffamily\bfseries\color{black}---\nobreakspace#3.}} % Optional theorem note
\renewcommand{\qedsymbol}{$\blacksquare$}% Optional qed square
\makeatother

% Defines the theorem text style for each type of theorem to one of the three styles above
\newcounter{dummy}
\numberwithin{dummy}{section}
\theoremstyle{ocrenumbox}
\newtheorem{theoremeT}[dummy]{Theorem}
\newtheorem{problem}{Problem}[chapter]
\newtheorem{exerciseT}{Exercise}[chapter]
\theoremstyle{blacknumex}
\newtheorem{exampleT}{Example}[chapter]
\theoremstyle{blacknumbox}
\newtheorem{vocabulary}{Vocabulary}[chapter]
\newtheorem{definitionT}{Definition}[section]
\newtheorem{corollaryT}[dummy]{Corollary}
\theoremstyle{ocrenum}
\newtheorem{proposition}[dummy]{Proposition}

%----------------------------------------------------------------------------------------
%	DEFINITION OF COLORED BOXES
%----------------------------------------------------------------------------------------

\RequirePackage[framemethod=default]{mdframed} % Required for creating the theorem, definition, exercise and corollary boxes

% Theorem box
\newmdenv[skipabove=7pt,
skipbelow=7pt,
backgroundcolor=black!5,
linecolor=ocre,
innerleftmargin=5pt,
innerrightmargin=5pt,
innertopmargin=5pt,
leftmargin=0cm,
rightmargin=0cm,
innerbottommargin=5pt]{tBox}

% Exercise box
\newmdenv[skipabove=7pt,
skipbelow=7pt,
rightline=false,
leftline=true,
topline=false,
bottomline=false,
backgroundcolor=ocre!10,
linecolor=ocre,
innerleftmargin=5pt,
innerrightmargin=5pt,
innertopmargin=5pt,
innerbottommargin=5pt,
leftmargin=0cm,
rightmargin=0cm,
linewidth=4pt]{eBox}

% Definition box
\newmdenv[skipabove=7pt,
skipbelow=7pt,
rightline=false,
leftline=true,
topline=false,
bottomline=false,
linecolor=ocre,
innerleftmargin=5pt,
innerrightmargin=5pt,
innertopmargin=0pt,
leftmargin=0cm,
rightmargin=0cm,
linewidth=4pt,
innerbottommargin=0pt]{dBox}

% Corollary box
\newmdenv[skipabove=7pt,
skipbelow=7pt,
rightline=false,
leftline=true,
topline=false,
bottomline=false,
linecolor=gray,
backgroundcolor=black!5,
innerleftmargin=5pt,
innerrightmargin=5pt,
innertopmargin=5pt,
leftmargin=0cm,
rightmargin=0cm,
linewidth=4pt,
innerbottommargin=5pt]{cBox}

% Creates an environment for each type of theorem and assigns it a theorem text style from the "Theorem Styles" section above and a colored box from above
\newenvironment{theorem}{\begin{tBox}\begin{theoremeT}}{\end{theoremeT}\end{tBox}}
\newenvironment{exercise}{\begin{eBox}\begin{exerciseT}}{\hfill{\color{ocre}\tiny\ensuremath{\blacksquare}}\end{exerciseT}\end{eBox}}
\newenvironment{definition}{\begin{dBox}\begin{definitionT}}{\end{definitionT}\end{dBox}}
\newenvironment{example}{\begin{exampleT}}{\hfill{\tiny\ensuremath{\blacksquare}}\end{exampleT}}
\newenvironment{corollary}{\begin{cBox}\begin{corollaryT}}{\end{corollaryT}\end{cBox}}

%----------------------------------------------------------------------------------------
%	REMARK ENVIRONMENT
%----------------------------------------------------------------------------------------

\newenvironment{remark}{\par\vspace{10pt}\small % Vertical white space above the remark and smaller font size
\begin{list}{}{
\leftmargin=35pt % Indentation on the left
\rightmargin=25pt}\item\ignorespaces % Indentation on the right
\makebox[-2.5pt]{\begin{tikzpicture}[overlay]
\node[draw=ocre!60,line width=1pt,circle,fill=ocre!25,font=\sffamily\bfseries,inner sep=2pt,outer sep=0pt] at (-15pt,0pt){\textcolor{ocre}{R}};\end{tikzpicture}} % Orange R in a circle
\advance\baselineskip -1pt}{\end{list}\vskip5pt} % Tighter line spacing and white space after remark

%----------------------------------------------------------------------------------------
%	SECTION NUMBERING IN THE MARGIN
%----------------------------------------------------------------------------------------

\makeatletter
\renewcommand{\@seccntformat}[1]{\llap{\textcolor{ocre}{\csname the#1\endcsname}\hspace{1em}}}
\renewcommand{\section}{\@startsection{section}{1}{\z@}
{-4ex \@plus -1ex \@minus -.4ex}
{1ex \@plus.2ex }
{\normalfont\large\sffamily\bfseries}}
\renewcommand{\subsection}{\@startsection {subsection}{2}{\z@}
{-3ex \@plus -0.1ex \@minus -.4ex}
{0.5ex \@plus.2ex }
{\normalfont\sffamily\bfseries}}
\renewcommand{\subsubsection}{\@startsection {subsubsection}{3}{\z@}
{-2ex \@plus -0.1ex \@minus -.2ex}
{.2ex \@plus.2ex }
{\normalfont\small\sffamily\bfseries}}
\renewcommand\paragraph{\@startsection{paragraph}{4}{\z@}
{-2ex \@plus-.2ex \@minus .2ex}
{.1ex}
{\normalfont\small\sffamily\bfseries}}

%----------------------------------------------------------------------------------------
%	PART HEADINGS
%----------------------------------------------------------------------------------------

% numbered part in the table of contents
\newcommand{\@mypartnumtocformat}[2]{%
\setlength\fboxsep{0pt}%
\noindent\colorbox{ocre!20}{\strut\parbox[c][.7cm]{\ecart}{\color{ocre!70}\Large\sffamily\bfseries\centering#1}}\hskip\esp\colorbox{ocre!40}{\strut\parbox[c][.7cm]{\linewidth-\ecart-\esp}{\Large\sffamily\centering#2}}}%
%%%%%%%%%%%%%%%%%%%%%%%%%%%%%%%%%%
% unnumbered part in the table of contents
\newcommand{\@myparttocformat}[1]{%
\setlength\fboxsep{0pt}%
\noindent\colorbox{ocre!40}{\strut\parbox[c][.7cm]{\linewidth}{\Large\sffamily\centering#1}}}%
%%%%%%%%%%%%%%%%%%%%%%%%%%%%%%%%%%
\newlength\esp
\setlength\esp{4pt}
\newlength\ecart
\setlength\ecart{1.2cm-\esp}
\newcommand{\thepartimage}{}%
\newcommand{\partimage}[1]{\renewcommand{\thepartimage}{#1}}%
\def\@part[#1]#2{%
\ifnum \c@secnumdepth >-2\relax%
\refstepcounter{part}%
\addcontentsline{toc}{part}{\texorpdfstring{\protect\@mypartnumtocformat{\thepart}{#1}}{\partname~\thepart\ ---\ #1}}
\else%
\addcontentsline{toc}{part}{\texorpdfstring{\protect\@myparttocformat{#1}}{#1}}%
\fi%
\startcontents%
\markboth{}{}%
{\thispagestyle{empty}%
\begin{tikzpicture}[remember picture,overlay]%
\node at (current page.north west){\begin{tikzpicture}[remember picture,overlay]%
\fill[ocre!20](0cm,0cm) rectangle (\paperwidth,-\paperheight);
\node[anchor=north] at (4cm,-3.25cm){\color{ocre!40}\fontsize{220}{100}\sffamily\bfseries\thepart};
\node[anchor=south east] at (\paperwidth-1cm,-\paperheight+1cm){\parbox[t][][t]{8.5cm}{
\printcontents{l}{0}{\setcounter{tocdepth}{1}}%
}};
\node[anchor=north east] at (\paperwidth-1.5cm,-3.25cm){\parbox[t][][t]{15cm}{\strut\raggedleft\color{white}\fontsize{30}{30}\sffamily\bfseries#2}};
\end{tikzpicture}};
\end{tikzpicture}}%
\@endpart}
\def\@spart#1{%
\startcontents%
\phantomsection
{\thispagestyle{empty}%
\begin{tikzpicture}[remember picture,overlay]%
\node at (current page.north west){\begin{tikzpicture}[remember picture,overlay]%
\fill[ocre!20](0cm,0cm) rectangle (\paperwidth,-\paperheight);
\node[anchor=north east] at (\paperwidth-1.5cm,-3.25cm){\parbox[t][][t]{15cm}{\strut\raggedleft\color{white}\fontsize{30}{30}\sffamily\bfseries#1}};
\end{tikzpicture}};
\end{tikzpicture}}
\addcontentsline{toc}{part}{\texorpdfstring{%
\setlength\fboxsep{0pt}%
\noindent\protect\colorbox{ocre!40}{\strut\protect\parbox[c][.7cm]{\linewidth}{\Large\sffamily\protect\centering #1\quad\mbox{}}}}{#1}}%
\@endpart}
\def\@endpart{\vfil\newpage
\if@twoside
\if@openright
\null
\thispagestyle{empty}%
\newpage
\fi
\fi
\if@tempswa
\twocolumn
\fi}

%----------------------------------------------------------------------------------------
%	CHAPTER HEADINGS
%----------------------------------------------------------------------------------------

% A switch to conditionally include a picture, implemented by  Christian Hupfer
\newif\ifusechapterimage
\usechapterimagetrue
\newcommand{\thechapterimage}{}%
\newcommand{\chapterimage}[1]{\ifusechapterimage\renewcommand{\thechapterimage}{#1}\fi}%
\newcommand{\autodot}{.}
\def\@makechapterhead#1{%
{\parindent \z@ \raggedright \normalfont
\ifnum \c@secnumdepth >\m@ne
\if@mainmatter
\begin{tikzpicture}[remember picture,overlay]
\node at (current page.north west)
{\begin{tikzpicture}[remember picture,overlay]
\node[anchor=north west,inner sep=0pt] at (0,0) {\ifusechapterimage\includegraphics[width=\paperwidth]{\thechapterimage}\fi};
\draw[anchor=west] (\Gm@lmargin,-9cm) node [line width=2pt,rounded corners=15pt,draw=ocre,fill=white,fill opacity=0.5,inner sep=15pt]{\strut\makebox[22cm]{}};
\draw[anchor=west] (\Gm@lmargin+.3cm,-9cm) node {\huge\sffamily\bfseries\color{black}\thechapter\autodot~#1\strut};
\end{tikzpicture}};
\end{tikzpicture}
\else
\begin{tikzpicture}[remember picture,overlay]
\node at (current page.north west)
{\begin{tikzpicture}[remember picture,overlay]
\node[anchor=north west,inner sep=0pt] at (0,0) {\ifusechapterimage\includegraphics[width=\paperwidth]{\thechapterimage}\fi};
\draw[anchor=west] (\Gm@lmargin,-9cm) node [line width=2pt,rounded corners=15pt,draw=ocre,fill=white,fill opacity=0.5,inner sep=15pt]{\strut\makebox[22cm]{}};
\draw[anchor=west] (\Gm@lmargin+.3cm,-9cm) node {\huge\sffamily\bfseries\color{black}#1\strut};
\end{tikzpicture}};
\end{tikzpicture}
\fi\fi\par\vspace*{270\p@}}}

%-------------------------------------------

\def\@makeschapterhead#1{%
\begin{tikzpicture}[remember picture,overlay]
\node at (current page.north west)
{\begin{tikzpicture}[remember picture,overlay]
\node[anchor=north west,inner sep=0pt] at (0,0) {\ifusechapterimage\includegraphics[width=\paperwidth]{\thechapterimage}\fi};
\draw[anchor=west] (\Gm@lmargin,-9cm) node [line width=2pt,rounded corners=15pt,draw=ocre,fill=white,fill opacity=0.5,inner sep=15pt]{\strut\makebox[22cm]{}};
\draw[anchor=west] (\Gm@lmargin+.3cm,-9cm) node {\huge\sffamily\bfseries\color{black}#1\strut};
\end{tikzpicture}};
\end{tikzpicture}
\par\vspace*{270\p@}}
\makeatother

%----------------------------------------------------------------------------------------
%	HYPERLINKS IN THE DOCUMENTS
%----------------------------------------------------------------------------------------

\usepackage{hyperref}
\hypersetup{hidelinks,backref=true,pagebackref=true,hyperindex=true,colorlinks=false,breaklinks=true,urlcolor=ocre,bookmarks=true,bookmarksdepth=3,bookmarksopen=false,pdftitle={Dev Web - CR TP},pdfauthor={Axel LE BOT}}
\usepackage{bookmark}
\bookmarksetup{
open,
numbered,
addtohook={%
\ifnum\bookmarkget{level}=0 % chapter
\bookmarksetup{bold}%
\fi
\ifnum\bookmarkget{level}=-1 % part
\bookmarksetup{color=ocre,bold}%
\fi
}
}
 % Insert the commands.tex file which contains the majority of the structure behind the template

\begin{document}

%----------------------------------------------------------------------------------------

%----------------------------------------------------------------------------------------
%	TITLE PAGE
%----------------------------------------------------------------------------------------

\begingroup
\thispagestyle{empty}
\begin{tikzpicture}[remember picture,overlay]
\node[inner sep=0pt] (background) at (current page.center) {\includegraphics[width=\paperwidth]{./Pictures/background.png}};
\draw (current page.center) node [fill=ocre!30!white,fill opacity=0.6,text opacity=1,inner sep=1cm]{\Huge\centering\bfseries\sffamily\parbox[c][][t]{\paperwidth}{\centering Module ITC315 - Web\\[15pt] % Book title
{\Large Compte-rendu de travaux pratiques}\\[20pt] % Subtitle
{\huge Basile DUTERTE | Léonard GÉHIN | Axel LE BOT | Andrew LENC}\\ % Author name
}};
\end{tikzpicture}
\vfill
\endgroup

%----------------------------------------------------------------------------------------
%	COPYRIGHT PAGE
%----------------------------------------------------------------------------------------

\newpage
~\vfill
\thispagestyle{empty}

\noindent Copyright \copyright\ 2017-2018 Axel LE BOT \\ % Copyright notice

\noindent Licensed under the Creative Commons Attribution-NonCommercial 3.0 Unported License (the ''License''). You may not use this file except in compliance with the License. You may obtain a copy of the License at \url{http://creativecommons.org/licenses/by-nc/3.0}. Unless required by applicable law or agreed to in writing, software distributed under the License is distributed on an \textsc{``as is'' basis, without warranties or conditions of any kind}, either express or implied. See the License for the specific language governing permissions and limitations under the License.\\ % License information

%----------------------------------------------------------------------------------------
%	INTRODUCTION
%----------------------------------------------------------------------------------------

%\pagestyle{empty} % No headers

%\chapterimage{./Pictures/cover-start}

%\chapter*{Introduction}\addcontentsline{toc}{part}{\texorpdfstring{\protect\@myparttocformat{Introduction}}{Introduction}}

%Ce rapport contient le compte-rendu des travaux pratiques de développement web (Module ITC34).\\
%En vous souhaitant une bonne lecture.

%\cleardoublepage % Forces the table of contents chapter to start on an odd page so it's on the right

%\pagestyle{fancy} % Print headers again

%----------------------------------------------------------------------------------------
%	TABLE OF CONTENTS
%----------------------------------------------------------------------------------------

%\usechapterimagefalse % If you don't want to include a chapter image, use this to toggle images off - it can be enabled later with \usechapterimagetrue

\chapterimage{./Pictures/cover-table_of_contents} % Table of contents heading image

\pagestyle{empty} % No headers

\tableofcontents % Print the table of contents itself

\cleardoublepage % Forces the table of content to start on an odd page so it's on the right

\pagestyle{fancy} % Print headers again

%----------------------------------------------------------------------------------------
%	TABLE OF FIGURES
%----------------------------------------------------------------------------------------

%\chapterimage{./Pictures/cover-table_of_figures} % Table of contents heading image

%\pagestyle{empty} % No headers

%\listoffigures % Print the table of contents itself

%\cleardoublepage % Forces the first chapter to start on an odd page so it's on the right

%\pagestyle{fancy} % Print headers again

%----------------------------------------------------------------------------------------
%	PART ONE
%----------------------------------------------------------------------------------------

%\part{TP}

%% exercice pair -> Leonard + Basile // impair -> Andrew + Axel

\chapterimage{./Pictures/cover-cobweb}
\chapter{TP1}
\section{Introduction}
Les sites internet sont téléchargés puis affichés à l'utilisateur par les navigateurs internet en interprétant les langages HTML et CSS.

\begin{itemize}
\item HTML signifie HyperText Markup Language : Permet de partitionner le contenu d'une page web.
\item CSS signifie Cascade Style Sheet : Permet d'agencer le contenu d'une page web.
\end{itemize}

Nous verrons lors de ce TP les utilisations basiques de ces langages.

%%%%%%%%%%%%%%%%%%%%%%%%%%%%%%%
%        Exercice 1
%%%%%%%%%%%%%%%%%%%%%%%%%%%%%%%
\section{Images fluides}
L'affichage d'image permet parfois de récapituler des idées, de les illustrer ou même simplement de décorer un site web. Chaque support dispose d'une taille d'écran unique, ainsi, afficher une image d'une taille fixée (en pixels) n'est pas une bonne idée si on veut que celle-ci s'affiche de manière similaire sur tout les appareils. Si une image de 800x600px apparaîta commeassew grande sur un écran de 1024x768px, elle apparaîtra gigantesque sur un terminal de 320x200px ou ridiculement petite sur un écran en 1440p. C'est pourquoi on choisit d'afficher les images en fonction d'un certain pourcentage de la largeur de l'écran. Dans l'exemple, on décide d'afficher le texte sur les 60 premiers pourcents de l'écran, puis d'afficher l'image sur les 40 derniers pourcents. Au premier essai, en mettant simplement:
\begin{minted}{css}
div.img-wrap{
	float:right;
	width:40%;
}
\end{minted}
L'image se met effectivement à droite du texte, mais elle dépasse largement de l'écran. On rajoute donc une classe "responsive", qu'on attribue a la balise img contenant l'image:
\begin{minted}{css}
.responsive{
  max-width: 100%;
  height: auto;
}
\end{minted}
Ainsi, l'image est limitée à 100% de la taille du div contenant l'image. Celui-ci étant limité à 40%, l'image sera donc redimensionnée pour tenir dans cette espace. On ajoute height:auto de façon redondante pour préciser la taille supposée de l'image.

%%%%%%%%%%%%%%%%%%%%%%%%%%%%%%%
%        Exercice 2
%%%%%%%%%%%%%%%%%%%%%%%%%%%%%%%
\section{Videos fluides}
Dans cet deuxième exercice du TP. Le but est d'obtenir un affichage interactif d'une vidéo sur une page web.\\
En HTML deux type de balises sont disponibles pour réaliser cet action.\\
En premier lieu nous avons la balise <video> qui permet d'afficher un fichier vidéo présent sur l'ordinateur sur la page web, et ceci depuis n'importe quelle répertoire.\\

La deuxième méthode consiste à ajouter une vidéo déjà présente sur internet et ceux grâce à son URL. Pour cela on utilise une balise différente nommé <iframe>. Cet balise recherche l'URL de la vidéo afin de redirigé l'affichage sur notre page web.\\

Dans le code on peut voir ici la balise iframe, ce qui indique l'ajout de la vidéo via le web.
\begin{minted}[linenos]{html}
<div class="video-wrap">
	<iframe src="https://www.youtube.com/embed/vI6as6rumJk" width="500" height="281" frameborder="0" autoplay="none"></iframe>
</div>
\end{minted}

Chaque vidéo peut être lu tel quel, cependant on peut aussi choisir de modifier la taille de sa fenêtre par exemple. C'est le cas ici. Pour cela on utilise deux instructions qui sont "width" et "height" , respectivement pour la longueur et la hauteur de la fenêtre en pixel. Il faut faire attention à ne pas négliger la taille réel de la vidéo.\\
Chaque instruction ou balise doit être placé dans une balise <div>. Celle ci permet de créer des divisions d'exécution. L'idée est d'ainsi de mettre une balise \mintinline{html}{<div>} autour des instructions de tailles également.

\begin{minted}[linenos]{html}
<video src="media/speaker.mp4" controls></video>
\end{minted}

Ici il s'agit du même principe, excepté le fait que l'ont utilise une balise <video> dans le sac.\\
Par ailleurs, on remarque l'utilisation dans les deux cas précédent d'une balise portant l'intitulé "contrôle".\\
Cette balise permet d'accéder aux fonctionnalités de base tel que la lecture  la pause et autre.\\

En ce qui concerne le css, on peut visuellement remarquer plusieurs parties.\\
Une première parti qui a pour but de gérer l'espace occupé par la vidéo. Il s'agit du \mintinline{css}{.video-wrap iframe} pour une vidéo web ou bien \mintinline{css}{.video-wrap video} si elle provient d'un espace de l'ordinateur. Ces balises ordonne la vidéo de telle sorte qu'elle prenne toute la place disponible.\\

Une fois la vidéo géré sur la page web, on peut géré la vidéo à l'intérieur de sa propre fenêtre. C'est la partie \mintinline{css}{.video-wrap} qui permet ceci. Son utilisation est fréquentes.

%%%%%%%%%%%%%%%%%%%%%%%%%%%%%%%
%        Exercice 3
%%%%%%%%%%%%%%%%%%%%%%%%%%%%%%%
\section{Redimension des images et requêtes media}
Les requêtes \mintinline{css}{media} permettent de conditionner l'affichage suivant les médias (écran, appareil, affichage ...).\\
Dans le fichier CSS nous pouvons ainsi spécifier plusieurs taille d'affichage :
\begin{itemize}
  \item Ecran taille <= 1024px
  \begin{minted}{css}
    @media screen and (max-width: 1024px) {
  \end{minted}
  \item Ecran 1025px <= taille <= 1280px
  \begin{minted}{css}
    @media screen and (min-width: 1025px) and (max-width: 1280px) {
  \end{minted}
  \item Ecran 1081px <= taille
  \begin{minted}{css}
    @media screen and (min-width: 1281px) {
  \end{minted}
\end{itemize}

%%%%%%%%%%%%%%%%%%%%%%%%%%%%%%%
%        Exercice 4 %Done
%%%%%%%%%%%%%%%%%%%%%%%%%%%%%%%
\section{Navigation adaptative avec des requetes media}
Dans cette partie, nous allons voir comment adapter efficacement la mise en page d'une page web. Ici, nous adapterons la page par rapport a la largeur de l'écran de l'utilisateur. Pour chaque catégorie de largeur d'écrans, on donnera des attribut différents aux éléments de la page.

Comme dans la partie precedente, nous allons utiliser des requetes media. Les ecrans des utilisateurs seront categorisés en trois categories en fonction de leurs largeurs:
\begin{itemize}
\item moins de 800 pixels
\item entre 801 et 1024 pixels
\item plus de 1025 pixels
\end{itemize}

Avec un écran très large, on peut faire plusieurs colonnes et mettre le menu sous forme de liste à puces  sur la gauche de la page.  La requête media permet de ne selectionner que les écrans de plus de 1024 pixels. Pour la liste à puces, le grand menu, on met son display en "left" et sa largeur a 20%. Le contenu occupe alors les 80% restant et est mis en "float right". Le "float" permet d'aligner un élément sur la droite, la gauche ou le milieu de son élément parent. Ici, les deux float differents permettent d'être sûr que le contenu et le menu soit bien alignés. Le "display:none" permet de cacher le petit menu, inutile ici.

Style utilisé pour les écrans larges:
\begin{minted}[linenos]{css}
@media screen and ( min-width: 1025px ) {
	.petit-menu { display: none; }
	.grand-menu {
		display: inline;
		float: left;
		width: 20%; /*Laisse 80% de place en largeur pour le texte*/
		}
	.contenu{
		float: right;
		width: 80%;
		}
}
\end{minted}

Pour un écran moins large, on utilise un affichage en liste à puces, mais cette fois, en ligne et en tête de page. Un tel affichage permet de gagner de la place en largeur. L'écran reste tout de même suffisemment large pour que le menu ne prenne pas plus d'une ligne.\\
Au niveau du code, on utilise une requête media pour etre utilise si l'écran a une largeur entre 801 et 1024 pixels. Comme sur l'écran large, le petit menu n'est pas affiché grâce a l'attribut "display". Pour le grand menu, qui doit être affiché en une ligne en en-tête, on met son display en inline et prend toute la largeur. On désactive l'attribut "list-style-type" de la balise <ul>, ce qui permet de manipuler des balises <li> comme des "blocs". Néanmoins, nous tranformont les balises <li> en "inline" pour qu'il n'y ait pas de retour à la ligne.

Style utilisé pour les écrans moyens:
\begin{minted}[linenos]{css}
@media screen and ( min-width: 801px ) and ( max-width: 1024px ) {
	.petit-menu { display: none; }
	.grand-menu {
		display:inline;
		width: 100%;
		}
	.grand-menu ul { list-style-type: none; }
	.grand-menu ul li { display: inline; }
	.contenu: { width: 100%; }
}
\end{minted}

Pour la plus petite catégorie d'écran, on utilise une liste déroulante placée au dessus du texte. C'est la solution la plus réduite pour placer un menu. Mais cette taille a un coût pour l'utilisateur : deux cliques / tapes du doigts seront necessaires pour utiliser ce menu.  Dans la feuille de style CSS, on applique simplement ce qu'il vient d'être vu : on cache le grand menu, le contenu prend toute la place qu'il peut et le petit menu est affiché. Mettre le petit menu en "block" ou "inline" n'influence pas sur le résultat final.

Style utilisé pour les petits écrans:
\begin{minted}[linenos]{css}
@media screen and ( max-width: 800px ) {
	.petit-menu { display: inline; }
	.grand-menu { display: none; }
	.contenu: { width: 100%; }
}
\end{minted}

%%%%%%%%%%%%%%%%%%%%%%%%%%%%%%%
%        Exercice 5
%%%%%%%%%%%%%%%%%%%%%%%%%%%%%%%
\section{Minimum et Maximum}
Afin de rendre le site plus compact et plus reactif, nous faisons floater nos trois élément à gauche et nous appliquons un taille maximum au deux dernier élément pour les afficher les un à la suite des autres sans que l'élément deux et trois empietent sur leur placement mutuel.

Des propriétés CSS nous permettent de borner la taille d'un element HTML :
\begin{itemize}
\item \mintinline{css}{min-width} et \mintinline{css}{max-width}
\item \mintinline{css}{min-height} et \mintinline{css}{max-height}
\end{itemize}

Nous souhaitons limiter uniquement la largeur de nos éléments HTML, ainsi nous ajoutons le code suivant aux classes deux et trois :
\begin{minted}{css}
max-width: 300px;
float: left;
\end{minted}

%%%%%%%%%%%%%%%%%%%%%%%%%%%%%%%
%        Exercice 6
%%%%%%%%%%%%%%%%%%%%%%%%%%%%%%%
\section{Controler la mise en page avec un padding relatif}
Dans une page web, il est parfois nécessaire et utile de placer des espaces vides entres ou dans certains éléments. Ces espaces peuvent être manipule via les attribut "margin" et "padding" en CSS. Le margin sert a ajouter un espace vide en dehors de la bordure de l’élément. Au contraire, le padding permet de placer un espace blancs entre le contenu de l’élément et sa bordure. Ces deux attribut peuvent avoir des longueur d’espaces blancs différentes dans chaque directions, généralement exprimées avec des valeurs s’adaptant a l’affichage de l’utilisateur (em, pourcentages, etc...).
\\\\
Dans la page web étudiée ici, nous allons utiliser le padding afin de rendre plus lisible un espace de commentaires. Ici, nous utiliserons la forme de padding ou nous explicitons les quatre valeur de padding, mais il est important de savoir que d’autre manière d’écrire le padding existent afin de factoriser la feuille de style. Lorsqu’on utilise l’écrire non factorisée, il peut être difficile de retenir l’ordre des valeur a mettre: haute, droite, bas, et gauche. Pour les retenir plus simplement, on peut se dire qu’on suit l’aiguille du montre.
\\\\
Dans un espace commentaire, les utilisateurs se répondent entre eux, des commentaires sont donc imbrique les uns dans les autres. Lorsque nous avons un affichage d’élément imbrique les uns dans les autres avec une hiérarchie, pouvoir immédiatement identifier les hiérarchies supérieures est très important pour les utilisateurs. Ici, nous utilisons le padding afin de mettre une marge a sur le cote gauche. Comme chaque commentaire possède une marge a sa gauche, tous les commentaire de hiérarchie inférieures seront décalés et ainsi de suite. Pour l’affichage des commentaires, nous avons utilise le padding ci-dessous. Vous pourez remarquer que nous avons exprime trois longueurs en pixels, ce qui est déconseille comme expliquer plus haut. Cependant, ces trois valeurs de padding ne servent qu’a écarter legerement les lignes des bordures des commentaires, nous pouvons donc nous permettre d’utiliser des pixel, unité de mesure fine.
\\\\
Padding utilisé pour les commentaires
\begin{minted}[linenos]{css}
.commentaire{
	padding: 8px 2px 2px 2em;
}
\end{minted}
Les attributs "padding" et "margin" permettent généralement d’espacer les éléments, et ainsi, rendent la page plus facile et agréable a lire pour l’utilisateur. Ces notions est donc très importante a maîtriser dans le développement web.

%%%%%%%%%%%%%%%%%%%%%%%%%%%%%%%
%        Exercice 7
%%%%%%%%%%%%%%%%%%%%%%%%%%%%%%%
\section{Tous terminaux}
L'accès à Internet est possible depuis une multitude d'appareils différents, il est ainsi parfois nécessaire de préciser un style différent pour certains appareils particuliers qui présente une interface inhabituelle ou moins pratique qu'une souris et un clavier. Par exemple, une télévision dont le seul moyen d'interaction est une télécommande qui répond lentement et n'est pas ergonomique, risque de décourager un utilisateur d'utiliser un site web. On va donc identifier le medium utilisé et s'adapter.
\begin{minted}{css}
@media (orientation: portrait){
  h1{
    width:100%;
    margin: auto;
    text-align:center;
  }
  img{
    clear:both;
    display:block;
    width:100%;
  }
}
@media (orientation: landscape){
  img{
    float:left;
  }
}
\end{minted}
Pour l'orientation portrait, on place le titre au milieu, mais on met aussi le titre sous l'image plutôt qu'à côté pour tirer parti de l'espace vertical. En paysage, on aligne les deux pour utiliser l'espace horizontale et donner un rendu plus naturel.
\begin{minted}{css}
@media screen and (min-width:1281px){
  body{
    margin: auto 10% auto 10%;
  }
  h1{
    width:100%;
    margin: auto;
    text-align:center;
  }
  img{
    float:left;
    display:block;
  }
}
\end{minted}
Pour les écrans de plus de 1280px de large, on centre le contenu en laissant une marge pour ne pas forcer l'utilisateur à trop bouger les yeux pour lire le contenu. On met aussi le texte et l'image sur une même ligne pour garder le défilement au minimum.
%%%%%%%%%%%%%%%%%%%%%%%%%%%%%%%
%        Conclusion
%%%%%%%%%%%%%%%%%%%%%%%%%%%%%%%
\section{Conclusion}
À travers ces 7 exercices, nous en avons découvert un peu sur la création de contenu réactif et adapté aux supports variés qui peuvent maintenant accéder a Internet. L'ajout de contenu multimédia comme les images ou les vidéos permet de bien accompagner le contenu textuel d'un site web, et il est important d'adapter ce contenu au support. L'adaptation de la taille des images et de la disposition des sites web permet de rendre le contenu plus accessible et beaucoup plus facile d'accès (comme sur les téléphone mobile par exemple).

% exercice impair -> Leonard + Basile // pair -> Andrew + Axel
\chapter{TP2}
\section{Introduction}
Dans le cas de l'approfondissement des connaissances sur le devellopement Web et ce qui y est rattaché, nous avons effectué un deuxième TP composé lui aussi d'une série d'exercice qui ont pour but de travailler des exemples utiles. La spécification propre à tous ces exercices est qu'il s'agit de devellopement Web pour à but d'utilisation sur les terminaux mobiles. Dans un premier temps, nous verrons comment adapter une page HTML suivant la taille du terminale, ceci en se basant sur divers exemple tel que l'utilation de formulaire mais aussi dans le cas de page d'information classique. Suite à ça, nous travaillerons sur la définition de fonctionnalité avec l'utilisation du java script, tout en prétant attention aux utilisations coté serveur/client. Après quoi nous verrons quelques bases sur le stockage local pour terminaux mobile.\\
Et pour finir, nous ferrons un point sur la géolocalisation qui est aujourd'hui présente sur tout les smartphones et utilisé par beaucoup d'application et de page web.

%%%%%%%%%%%%%%%%%%%%%%%%%%%%%%%
%        Exercice 1
%%%%%%%%%%%%%%%%%%%%%%%%%%%%%%%
\section{Etiquette Formulaire}
Après ouverture du fichier html, je créer un fichier css. On ajoute ensuite les quelques lignes précisé dans la diapo.
\begin{itemize}
\item \mintinline{css}{h1 {font-size: 18pt;}} ici, il s'agit uniquement de définir une taille et un type de police pour les balise "h1"
\item \mintinline{css}{width: 100px;}, cela correspond à une largeur de 100 pixels pour le blocs d'instruction dans le quel est compris le "width"
\item \mintinline{css}{text-align: right;} cette instruction force l'allignement du texte par la droite.
\item \mintinline{css}{display: inline-block;} Ici le display inline block permet une boite de bloc principale ainsi que plusieurs boîtes en ligne en suivant.
\item \mintinline{css}{vertival-align: baseline} pour ce cas, il s'agit uniquement de l'alignement vertical des blocs qui suivent.
\end{itemize}

Afin de permettre de changer l'affichage sur les écrans de moins de 480px, nous ajoutons quelque ligne sur le css.
En premier lieu, avec l'instruction @media screen, on précise que les instructons qui vont suivre s'applique uniquement pour les média qui dispose d'une écran.
En précisant le \mintinline{css}{and (max-width: 480px)} on affine notre choix on précisant les médias qui dispose d'un écran de taille max = 480px.
Suite à ça, la ligne \mintinline{css}{.col-25, .col-75} permet une répartition en pourcentage de l'espace qui sera occupé sur l'écran. À raison de 25% et 75% ici.
Les deux lignes qui suivent, sont une précision concernant la largeur maximun à utilisé ainsi que "margin-top" qui affecte un espace nul avec la marge du haut.

Après avoir affiché le fichier HTML dans le navigateur, on remarque une changement sur la disposition vue au préalable. En effet suivant la taille de la fenêtre ouverte, on remarque que les étiquettes sont affichées ou non au dessus des champs du formulaire. Ceci est dû aux commande du css rajouté plus tôt.


%%%%%%%%%%%%%%%%%%%%%%%%%%%%%%%
%        Exercice 2
%%%%%%%%%%%%%%%%%%%%%%%%%%%%%%%
\section{Media queries} %2
Afin de changer l'affichage du menu en une colonne à gauche, qui prend de la place de façon horizontale, en une disposition horizontale au dessus du contenu. La limite pour ce choix sera de 800 pixels.
On ajoute donc ce code, qu'on détaillera après:
 Pour les écrans supérieur ou égal à 801px
\begin{minted}[linenos]{css}
@media screen and (min-width: 801px) {
  #titrePage {
    font-size:36pt;
    font-family:"Times New Roman", Times, serif;
    background-color:#9999FF
  }
  #container {
    font-size: 16pt;
    position: relative;
    width: 100%;
  }

  #colonneG {
    width: 200px;
    height: 100%;
    float:left;
  }

  #contenuPage {
    margin-left: 210px;
  }

  #footer {
    border: 2px gray solid;
    padding: 5pt;
    margin-top: 5pt;
  }
}
\end{minted}
La partie correspondant aux grands écrans est fourni, il n'est donc pas nécessaire de la détailler. Celle-ci s'assure juste de créer une colonne à gauche avec le menu, pour faciliter l'accès à celui-ci. On laisse une marge sur le contenu pour l'aérer, et profiter de l'espace.\\
Pour les écrans inférieur ou égal à 800px

\begin{minted}[linenos]{css}
@media screen and (max-width: 800px) {
  #colonneG{
    display:inline-block;
    clear:both;
    }
  #colonneG ul li {
    float:left;
    list-style-type:none;
    margin:5px;
  }
  \end{minted}
   Afin de replacer la colonne de gauche au dessus, on utilise clear:both pour empêcher d'autres éléments de s'afficher à droite et à gauche. On précise float:left sur les éléments de la liste non-ordonnée afin de les afficher en ligne, cela déplace aussi le titre de la colonne, on ajoute display:inline-block afin d'éviter ce phénomène.
  \begin{minted}[linenos]{css}
  #colonneG ul {
    padding:0;
  }
  #titrePage {
    font-size:36pt;
    font-family:"Times New Roman", Times, serif;
    background-color:#e3c773;
    color:#e8e8e8;
    text-align:right;
  }
}
\end{minted}


 On aligne le menu à gauche avec padding:0, pour que celui-ci ne gaspille pas la place à gauche. Les puces en face de chaque entrée restent en dépit de leur alignement horizontal, on utilise donc liste-style-type:none afin de les retirer et on rajoute une marge de 5px pour éviter que les éléments apparaissent trop proches les uns des autres. Enfin, on ajuste le titre de la page pour que celui-ci s'aligne à gauche, et soit en gris-blanc sur un fond d'une autre couleur.

%%%%%%%%%%%%%%%%%%%%%%%%%%%%%%%
%        Exercice 3
%%%%%%%%%%%%%%%%%%%%%%%%%%%%%%%
\section{Détection de fonctionnalités}
Après avoir récupéré les fichiers disponibles, j'ajoute le chemin vers la librairie Modernizr ainsi que la balise \mintinline{html}{<html class="no-js">}. Ceci a pour but de pouvoir utiliser les options java script qui seront plus tard utilisée.
Je créer et inclue un nouveau fichier java script que je nomme \mintinline{html}{detection_de_fonctionnalité.js}.
Dans ce java script on ajoute plusieur ligne avec l'instruction \mintinline{js}{document.querySelector} Cet méthode fonctionne de manière simple. Elle retourne le premier element dans le document en question qui correspont au selecteur passé en paramètre. Si rien n'est trouvée, la méthode renvois null.
Le code en question ( mettre un screen, c'est la première question) test les différentes fonctionnalités lié a la méthode "document.querySelector". Ici on test plus particulièrement la prise en charge de geoloc,touch, svg et canvas. Par rapport a ça, un message est renvoyé qui indique l'état sur la prise en charge.
Suite à ça, on ajoute un fichier css, que l'on nomme \mintinline{html}{detection_de_fonctionnalité.css} dans le quel on implemente le code de l'énoncé.
Ce code fonctionne de manière très simple. Prenons une ligne pour exemple.
\mintinline{css}{.animtest, .noanimtest {display: none;}} Ici nous avons les instructions .animtest et .noanitest qui sont des futurs paramètre pour la méthode display. Ensuite il y a la méthode display qui en fonction de l'instruction mis en paramètre derrière aura différent rôle sur les instructions précédentes.
Les 3 lignes fonctionnent de manière identique.
Le \mintinline{css}{display: none;} permet de rendre visible tout ce que est concerné par l'affichage avant le display.
Et ici, le display block permet de forcer un élément à générer une boîte de bloc principale.

%%%%%%%%%%%%%%%%%%%%%%%%%%%%%%%
%        Exercice 4
%%%%%%%%%%%%%%%%%%%%%%%%%%%%%%%
\section{Stockage Local} %4

Afin de réutiliser les clés de valeur, nous stockons les valeurs dans les clées suivantes :
\begin{minted}{js}
var KEY_FIRST_NAME = "firstName";
var KEY_LAST_NAME = "lastName";
var KEY_POST_CODE = "postCode";
\end{minted}

\begin{minted}{js}
if (window.localStorage.length !== 0) {
        console.log("Input replace");
        document.getElementById(KEY_FIRST_NAME).value = window.localStorage.getItem(KEY_FIRST_NAME);
        document.getElementById(KEY_LAST_NAME).value = window.localStorage.getItem(KEY_LAST_NAME);
        document.getElementById(KEY_POST_CODE).value = window.localStorage.getItem(KEY_POST_CODE);
    }
\end{minted}

Au moment où nous appuyons sur le bouton `Enregistrer` du formulaire la function \mintinline{js}{function storeLocalContent(fName, lName, pCode)} sera appeler comme le stipule la ligne html suivante :
\begin{minted}{html}
<input type="button" value="Enregistrer" onclick="storeLocalContent(
                   document.querySelector('#firstName').value,
                   document.querySelector('#lastName').value,
                   document.querySelector('#postCode').value
               )">
\end{minted}


Afin de vider le stockage local nous ajoutons la fonction `clearLocalContent` qui permettra de supprimer les données rajouter dans le stockage local.
Pour ne pas interférer avec d'autre script qui peuvent aussi utiliser le stockage local, nous n'utiliserons pas la function `clear()` mais nous supprimerons tous les éléments clés utilisé uniquement dans ce script.

\begin{minted}[linenos]{js}
function clearLocalContent() {
    //window.localStorage.clear(); // Suppression des autres clés utilisés par d'autres scripts donc attention
    window.localStorage.removeItem(KEY_FIRST_NAME);
    window.localStorage.removeItem(KEY_LAST_NAME);
    window.localStorage.removeItem(KEY_POST_CODE);
}
\end{minted}
Cette fonction est appelé quand nous appuyons sur le bouton "Effacer" comme nous le stipule la ligne HTML ci-dessous :

Afin de lancer le script d'initialisation nous pouvons soit ajouter le lien vers le script à la fin du body, soit en remplacant la ligne \mintinline{js}{window.onload=} par :
\begin{minted}{js}
document.addEventListener('DOMContentLoaded', function() {
    initialize();
});
\end{minted}

Ainsi en ayant enregistrer le formulaire nous pouvons au rechargement de la page voir se ré-afficher les données du formulaire.

%%%%%%%%%%%%%%%%%%%%%%%%
%Exo 5: Geolocalisation
%%%%%%%%%%%%%%%%%%%%%%%%
\section{Géolocalisation}
Dans ce dernier exercice nous allons utiliser et manipuler un peu plus en détail la librairie "Modernizr.js". Comme dit dans le titre, nous allons principalement utiliser les fonctionnalités de geolocalisation et afficher certaines informations sur la page html.

Comme vu dans l'exercice trois de ce TP, il est necessaire que l'utilisateur autorise l'utilisation de la géolocalisation. Dans le cas ou la géolocalisation ne fonctionne pas, la fonction \mintinline{js}{geoError()} affiche à l'utilisateur une boite de dialogue lui indiquant pourquoi la géolocalisation ne fonctionne pas. Ci-dessous nous pouvons voir la fonction \mintinline{js}{geoSuccess()}, qui donne des valeurs à des balises span du code HTML de la base pour afficher les informations de géolocalisation. La variable \mintinline{js}{positionInfo} est une structure de donnée, tout comme les parametre de la fonction \mintinline{js}{calculDistance()}, à l'avant derniere ligne de code. Cette initialisation de variable, très peu utilisé dans les langage plus courrant tel que le Java ou le C, permet de ne pas sauvegarder la variable en mémoire , puisque utiliser une seule fois dans le code. Pour accéder aux variables contenu dans la structure, on pourra par exemple écrire \mintinline{js}{nomStruc.latitude}, où nomStruct est le nom du paramètre dans la fonction \mintinline{js}{calculDistance()}.

\begin{minted}[linenos]{js}
function geoSuccess(positionInfo) {
	document.getElementById("longitude").innerHTML = positionInfo.coords.longitude;
	document.getElementById("latitude").innerHTML = positionInfo.coords.latitude;
	document.getElementById("precision").innerHTML = positionInfo.coords.accuracy;
	document.getElementById("altitude").innerHTML = positionInfo.coords.altitude;
	document.getElementById("precisionAltitude").innerHTML = positionInfo.coords.altitudeAccuracy;
	document.getElementById("cap").innerHTML = positionInfo.coords.heading;
	document.getElementById("vitesse").innerHTML = positionInfo.coords.speed;
	document.getElementById("distance").innerHTML = calculDistance(
			{latitude : positionInfo.coords.latitude ,
			 longitude : positionInfo.coords.longitude },
			{latitude : 47.3121519, /*Position de l'esirem*/
			 longitude : 5.0039326 });
}
\end{minted}

Les deux fonctions précédemment vu sont passé en paramètre à une seule fonction : \mintinline{js}{getLocation()}. qui les appellera si il y a succes ou si il y a erreur. Ce sont des "callback".
\begin{itemize}
  \item Si la géolocalisation est pris en charge, on utilise la fonction \mintinline{js}{geoSuccess()}.
  \item Si la géolocalisation n'est pas prise en charge on utilise la fonction \mintinline{js}{geoError()}.
\end{itemize}

\begin{minted}[linenos]{js}
function getLocation() {
	if (Modernizr.geolocation) { //Envoie true si geoLocation est supporte
		navigator.geolocation.getCurrentPosition(geoSuccess, geoError);
	}
}
\end{minted}

La fonction \mintinline{js}{getLocation()} est appelée au chargement de la page à l'aide de l'affectation de \mintinline{js}{windows.onload}.
\begin{minted}[linenos]{js}
window.onload = getLocation();
\end{minted}


%%%%%%%%%%%%%%%%%%%%%%%%
%Conclusion
%%%%%%%%%%%%%%%%%%%%%%%%
\section{Conclusion}
À travers ces 5 exercices, nous en avons découvert un peu sur la création de contenu réactif et adapté aux supports variés qui peuvent maintenant accéder a Internet. L'ajout de contenu multimédia comme les images ou les vidéos permet de bien accompagner le contenu textuel d'un site web, tandis que l'utilisation des informations géographiques peut permettre d'offrir un contenu plus personnalisé si le site offre des services locaux.
Dans ces travaux pratiques, nous avons pus apprendre à utiliser le javascript et récupérer quelques données, tel que la position de l'utilisateur. Avec le javascript, il est possible de modifier le contenu de la page de facon dynamique, et sans retour serveur, le javascript est par conséquent un langage important pour cerner les bases du dévelopement web. Avec ce TP, nous avons pus constater que le javascript est un langage ne compilant pas mais s'executant pas à pas, il est donc possible de faire une faute de frappe et de se rendre compte que nous avons fait cette erreur-ci seulement trente minutes plus tard.
Nous avons pu aussi aborder le concept important du javascript qu'est le "callback" énomément utilisé par exemple sur des serveurs utilisant des plateformes comme NodeJS.


%----------------------------------------------------------------------------------------
%	CONCLUSION
%----------------------------------------------------------------------------------------

%\cleardoublepage % Forces the conclusion chapter to start on an odd page so it's on the right

%\pagestyle{empty}

%\chapterimage{./Pictures/cover-end}

%\chapter*{Conclusion}\addcontentsline{toc}{part}{\texorpdfstring{\protect\@myparttocformat{Conclusion}}{Conclusion}}

%\lipsum[3-4]

%\pagestyle{fancy}

%----------------------------------------------------------------------------------------

\end{document}
